\section{Алгоритмы комплексной навигации}

\subsection{Проблема шума и фильтрация}
Данные, получаемые от системы компьютерного зрения, всегда содержат шум. Координаты дрона <<дрожат>> даже если он висит неподвижно. Это происходит из-за дискретности матрицы камеры, изменений освещения и вибраций корпуса.

Если передавать эти <<сырые>> координаты напрямую в полетный контроллер, дрон будет дергаться, пытаясь отработать каждое ложное смещение.

\subsubsection{Фильтр Калмана (Kalman Filter)}
Это стандарт де-факто в робототехнике. Фильтр Калмана не просто усредняет значения, он предсказывает, где дрон \textit{должен} оказаться в следующий момент времени, и корректирует это предсказание на основе новых измерений.

Алгоритм работает циклически в два этапа:
\begin{enumerate}
    \item \textbf{Prediction (Предсказание):} На основе физической модели (скорость, ускорение) мы оцениваем новую позицию. При этом растет неопределенность (мы не уверены в прогнозе).
    \item \textbf{Update (Коррекция):} Мы получаем измерение от камеры. Сравниваем его с прогнозом. Чем точнее датчик (меньше шум измерения $R$), тем больше мы верим ему. Чем точнее наша модель (меньше шум процесса $Q$), тем больше верим прогнозу.
\end{enumerate}

\begin{codefiletitle}[python]{Пример 1D фильтра Калмана}{python/examples/advanced_navigation/kalman_1d_example.py}
\end{codefiletitle}

Выше приведен только ключевой метод \texttt{update()}. Полный запускаемый скрипт с классом \texttt{KalmanFilter1D} и модельным экспериментом находится в файле \texttt{advanced\_navigation/kalman\_filter.py} --- именно его нужно запускать в заданиях ниже.

\subsection{Системы координат}
Одна из самых частых ошибок новичков — путаница в системах координат.

\begin{itemize}
    \item \textbf{Body Frame (Связанная система):} Ось X всегда смотрит вперед (нос дрона), Y — влево/вправо, Z — вверх. Если дрон поворачивается, оси поворачиваются вместе с ним.
    \item \textbf{World Frame / NED (North-East-Down):} Ось X смотрит на Север, Y — на Восток, Z — Вниз (к земле). Эта система неподвижна.
    \item \textbf{Marker Frame:} Ось Z обычно смотрит перпендикулярно плоскости маркера (вверх от него).
\end{itemize}

Для навигации нам нужно переводить координаты маркера (полученные камерой в Body Frame) в глобальные координаты дрона (World Frame).
$$ Pos_{world} = Pos_{drone} + R_{drone} \cdot Pos_{marker\_relative} $$

\subsection{Точная посадка с PID-регулятором}
Для автоматической посадки на зарядную станцию используется PID-регулятор.
\begin{itemize}
    \item \textbf{P (Proportional):} Чем дальше мы от центра, тем сильнее наклоняемся к нему.
    \item \textbf{I (Integral):} Если нас сносит ветром, мы накапливаем ошибку и наклоняемся сильнее, чтобы компенсировать снос.
    \item \textbf{D (Derivative):} Если мы летим к центру слишком быстро, мы <<тормозим>> заранее, чтобы не пролететь мимо.
\end{itemize}

\begin{taskbox}[Задание 5: Сравнение фильтров]
Запустите скрипт \texttt{kalman\_filter.py}.
1. Измените \texttt{measurement\_noise} (шум измерения) с 0.1 до 5.0. Как ведет себя фильтр?
2. Сравните работу фильтра Калмана и простого скользящего среднего (Moving Average) при резком скачке координат (например, дрон толкнули). Кто быстрее реагирует?
\end{taskbox}

\begin{taskbox}[Задание 6: Посадка в круг]
Напишите скрипт, который:
1. Ищет маркер с высоты 1.5 метра.
2. Снижается до 0.5 метра, удерживая маркер в центре кадра (используя P-регулятор по скоростям $V_x, V_y$).
3. Если ошибка позиционирования $< 5$ см в течение 2 секунд — совершает посадку.
\end{taskbox}

\begin{questionbox}[Контрольные вопросы]
\begin{enumerate}
    \item Зачем нужен фильтр Калмана, если можно просто взять среднее арифметическое за 10 кадров?
    \item Что произойдет с PID-регулятором, если перепутать знак у коэффициента P?
    \item В какой системе координат (Body или World) камера выдает вектор \texttt{tvec}?
\end{enumerate}
\end{questionbox}
